\section{Grundlagen}

Im Folgenden werden Begriffe zu den Themen Sicherheit und Shift-Left in der Softwareentwicklung erklärt, um somit eine grundlegende Wissensbasis für die danach folgenden Kapitel zu schaffen.

Als \textit{sichere Software} definieren wir solche, die den grundlegenden Prinzipien der Informationssicherheit \cite{blakley_information_2001} entspricht und auch trotz verschiedener Widrigkeiten in ihrer Funktionalität zuverlässig bleibt \cite{oueslati_literature_2015}. Diese Prinzipien werden in folgende Kategorien unterteilt:
\begin{enumerate}
    \item \textbf{Vertraulichkeit} beschreibt den Schutz sensibler Daten vor unbefugtem Zugriff.
    \item \textbf{Integrität} stellt sicher, dass die Daten konsistent und unverändert bleiben, außer sie werden bewusst durch autorisierte Nutzer verändert.
    \item \textbf{Verfügbarkeit} garantiert, dass die Software und ihre genutzten Daten und Ressourcen stets zugänglich sind. Die Verfügbarkeit muss auch unter gezielten Angriffen gewährleistet werden.
    \item \textbf{Zuverlässigkeit} leistet ein System dann, wenn es sich zu jedem Zeitpunkt erwartungsgemäß verhält. Dies gilt ebenfalls unter etwaigen Widrigkeiten, wie beispielsweise unbefugten Zugriffen von außen.
\end{enumerate}
Diese vier Prinzipien bilden die Grundlage für sichere Software und müssen somit bei der Entwicklung von sicherheitskritischen Prozessen gezielt beachtet werden.\

Der \textit{Shift-Left} Ansatz beschreibt das Verschieben von Prozeduren bzw. Aktivitäten aus späteren Phasen des Software Development Life Cycle in Frühere \cite{andriadi_impact_2023}. An einer zeitlichen Achse, wie sie in der Abbildung \ref{fig:shiftleft} zu sehen ist, lässt sich die Namensherkunft des Shift-Left Ansatzes erklären. Die Abbildung zeigt, wie der Fokus der Softwarequalität von einem späteren Zeitpunkt in einen Früheren verlagert wird. An der Achse entlang gesehen, geschieht diese Verschiebung (engl. Shift) also nach links (engl. left). Diese Idee ist schon länger bekannt, doch der Name Shift-Left für diesen Prozess wurde durch \citet{smith_shift-left_2001} geprägt. Dieser sprach 2001 vom \textit{Shift-Left Testing} und erklärte, dass durch das frühere Testen in der Softwareentwicklung die Qualität der Software gesteigert und die Kosten von Fehlern gesenkt werden könne \cite{dawoud_better_2024}. Dafür müssten Quality Assurance (QA) und Entwickler parallel arbeiten, was zur Folge hätte, dass das QA-Team 1. nicht auf den gesamten Entwicklungsprozess warten muss und 2. entdeckte Fehler bereits während der Entwicklung direkt wieder beseitigt werden können \cite{andriadi_impact_2023}. Das Ziel dabei ist es Fehler früh zu Erkennen und somit späteren Kosten- oder Zeitaufwand zu sparen \cite{dawoud_better_2024}. Auch wenn Shift-Left anfangs für das Testing vorgesehen wurde, lässt das Konzept sich jedoch auf viele erdenkliche Bereiche anwenden.

\begin{figure}
    \centering
    \includegraphics[width=0.9\linewidth]{images/Shift_Left.png}
    \caption{Shift-Left Verschiebung von Fokus auf Qualität im Software Development Life Cycle \cite{andriadi_impact_2023}}
    \label{fig:shiftleft}
\end{figure}

Mit \textit{Shift-Left Security} ist die Anwendung des Shift-Left Gedanken auf den Bereich der Software-Security gemeint \cite{dawoud_better_2024}. Viele Sicherheitsprozesse werden in der Regel erst am Ende der Entwicklung durchgeführt, wodurch Fehler oder Sicherheitslücken oft zu spät erkannt werden \cite{dawoud_better_2024}. Durch Shift-Left Security werden jene Prozesse in die früheren Phasen des Software Development Cycles verschoben, wodurch Bedrohungen zeitnah erkannt und beseitigt werden können. Jedoch wurde in den eingangs genannten Problemen gezeigt, dass das Einbinden von Sicherheitsprozessen in das agile Arbeiten schwierig umzusetzen ist. Dafür werden im folgenden Konzepte und Methoden definiert, die bei der Umsetzung in agilen Umgebungen helfen. Die Wirksamkeit und genauere Richtlinien dieser Ideen, werden im Hauptteil erläutert.

\textit{Development, Security and Operations (DevSecOps)} beschreibt die Erweiterung des klassischen \textit{DevOps} durch die Einbindung von Security Aspekten \cite{rajapakse_challenges_2022}. Die Aufgabe des DevSecOps ist es, Sicherheitspraktiken in den gesamten Entwicklungszyklus einzubinden, um so die Sicherheit der Software zu gewährleisten. Dabei werden automatisierte Sicherheitstest kontinuierlich angepasst und ausgewertet, wodurch die Sicherheit gewährleistet wird, sowie auch Zeit und Kosten gespart werden \cite{lombardi_devops_2023}. DevSecOps zielt ebenfalls darauf ab, die Zusammenarbeit und Kommunikation von Entwicklern, Sicherheitsexperten und Betriebsteams zu verbessern, um ein Verständnis von Sicherheit im ganzen Team zu bilden \cite{rajapakse_challenges_2022}.


